% Generated draft: requires mathematical and source review.
\documentclass[11pt]{article}
\usepackage[margin=25mm]{geometry}
\usepackage{fontspec}
\setmainfont{Latin Modern Roman}
\usepackage{amsmath,amssymb,mathtools,unicode-math}
\setmathfont{Latin Modern Math}
\usepackage{enumitem,needspace,xurl,hyperref}
\setlength{\parindent}{0pt}
\setlength{\parskip}{5pt}
\setlength{\emergencystretch}{3em}
\newcommand{\N}{\mathbb{N}}
\newcommand{\Z}{\mathbb{Z}}
\newcommand{\Q}{\mathbb{Q}}
\newcommand{\R}{\mathbb{R}}
\newcommand{\C}{\mathbb{C}}
\newcommand{\F}{\mathbb{F}}
\newcommand{\A}{\mathbb{A}}
\newcommand{\PP}{\mathbb{P}}
\newcommand{\E}{\mathbb{E}}
\newcommand{\eps}{\varepsilon}
\newcommand{\dd}{\,\mathrm{d}}
\begin{document}
\section*{UnsolvedMath ID: 30000943}
\textbf{Source problem code:} catalogue entry
\subsection{Definitions and mathematical statement}
In a standard alloy-type setting, a random Schrödinger operator has the form
\[
H_\omega=-\Delta+V_{\mathrm{per}}+\sum_{j\in\Z^d}\omega_j u(\,\cdot-j),
\]
with a sign-changing single-site potential $u$ and coupling law $\mu$ supported on $[a,b]$. The exact dimension, independence and potential hypotheses of the model considered in the report are not supplied; the display is contextual, not a replacement by all operators of this form.

Let $E_0$ be the bottom of the almost-sure spectrum and let $N(E)$ be the integrated density of states, defined through the appropriate normalized finite-volume eigenvalue-counting limit. Let $E_-(\lambda)$ be the lowest Neumann eigenvalue of the source's reference problem at coupling $\lambda$. Its cell, boundary geometry and potential are not specified in the extract.

The conjecture asks for Lifshitz-tail behavior of $N(E)$ as $E\downarrow E_0$ except possibly in the joint exceptional case
\[
\mu\text{ is Bernoulli}\quad\text{and}\quad E_-(a)=E_-(b).
\]
Equivalently, retain the exact sufficient-condition logic
\[
\bigl[\mu\text{ is not Bernoulli}\ \text{or}\ E_-(a)\ne E_-(b)\bigr]
\quad\Longrightarrow\quad
\text{Lifshitz tails at }E_0.
\]
The statement does not assert that both separate conditions are necessary, nor does ``unless'' alone assert failure of Lifshitz tails in the excluded joint case. In the usual nondegenerate endpoint-Bernoulli convention, $\mu=t\delta_a+(1-t)\delta_b$ with $a<b$ and $0<t<1$.

Lifshitz-tail terminology requires a quantitative lower-edge decay definition, for example a specified logarithmic exponent involving $\log|\log N(E)|/\log(E-E_0)$. The extract does not supply that definition, the exponent, endpoint probability bounds or the operator conditions needed to determine it. In particular, no exact exponent is imposed for an arbitrary coupling law. It is also unclear whether $a,b$ must be the actual extremal support points rather than merely endpoints of a containing interval. These qualifications must be recovered to obtain a fully specified analytic assertion.
\subsection{Short English statement}
For the source's sign-changing alloy model, ask for Lifshitz tails except possibly when the coupling law is Bernoulli and its endpoint Neumann energies coincide. The exception is conjunctive, and the model and precise tail notion require recovery.
\subsection{Sources}
\begin{itemize}
\item Supplied report: \emph{Random Schrödinger Operators and Random Matrices} (2009), Oberwolfach Reports, DOI 10.4171/owr/2008/16, pp. 797--868; proposed year 2008. \url{https://doi.org/10.4171/owr/2008/16}.
\item Supplied reference: Najar, \emph{Lifshitz tails for non-monotone models}, \url{https://www.numdam.org/item/SEDP_2007-2008____A14_0.pdf}.
\item Local UnsolvedMath registry, permanent ID 30000943, catalogue entry; created 25 August 2026, updated 16 September 2026. Its corrected\_verified note says the official report confirms the joint Bernoulli/equal-Neumann-energy exception. The supplied 21 August 2026 literature triage labels the general nonsymmetric assertion open, reporting later special cases rather than the full claim. No PDF access or independent current-status check is asserted here.
\end{itemize}
\end{document}
